% Methods section: construction of mock stellar halos with StreamHalo.
% AASTeX (aastex631) conventions; paste into the main manuscript.

\section{Mock Stellar Halos} \label{sec:mock}

We construct mock stellar halos by populating a Milky Way--like host potential with a
statistical sample of accreted satellites, disrupting each satellite into a tidal stream with
a particle-spray model, and embedding the resulting streams in a smooth, phase-mixed
background. The procedure is implemented in the \texttt{StreamHalo} package, which uses
the JAX-accelerated particle-spray code \texttt{StreaMAX} to generate the streams. Each step
is described below, and the fiducial parameters are listed in Table~\ref{tab:mock_params}.

\subsection{Host Potential} \label{sec:mock_host}

The host galaxy is a static, axisymmetric potential with four components: a
\citet{Miyamoto1975} stellar disk, \citet{Hernquist1990} bulge and nucleus components, and a
spherical \citet*[NFW;][]{Navarro1997} dark matter halo. All components share a common center
and their symmetry axis lies along $\hat{z}$. The component masses and scale lengths
(Table~\ref{tab:mock_params}) follow a standard Milky Way mass model. The same potential is
used to assign satellite velocities (Section~\ref{sec:mock_phase}), to compute tidal radii
(Section~\ref{sec:mock_streams}), and to integrate orbits.

\subsection{Satellite Population} \label{sec:mock_sats}

Satellite (sub)halo masses $M_h$ are drawn independently from a power-law mass function,
\begin{equation}
    \frac{dN}{dM_h} \propto M_h^{\,\alpha}, \qquad M_{\rm min} \le M_h \le M_{\rm max},
    \label{eq:hmf}
\end{equation}
by inverse-transform sampling, with a fiducial slope $\alpha = -1.9$, consistent with the
subhalo mass functions found in cosmological simulations \citep[e.g.,][]{Springel2008}.
Each halo mass is converted to a stellar mass $M_\star$ with the $z=0$ stellar-to-halo mass
relation of \citet[][their Equation~3]{Behroozi2013},
\begin{equation}
    \log_{10} M_\star = \log_{10}(\epsilon M_1) + f\!\left[\log_{10}(M_h/M_1)\right] - f(0),
\end{equation}
\begin{equation}
    f(x) = -\log_{10}\!\left(10^{\alpha_B x} + 1\right)
           + \delta\,\frac{\left[\log_{10}(1 + e^{x})\right]^{\gamma}}{1 + \exp(10^{-x})},
\end{equation}
with $\log_{10}(M_1/M_\odot) = 11.514$, $\log_{10}\epsilon = -1.777$, $\alpha_B = -1.412$,
$\delta = 3.508$, and $\gamma = 0.316$.

The number of satellites is not fixed in advance. We draw satellites one at a time until
their cumulative stellar mass first exceeds a target $M_{\star,\rm tot}$, which sets the
total accreted stellar mass of the halo. With the fiducial parameters
($M_{\star,\rm tot} = 10^{9}\,M_\odot$), a typical realization contains
$\sim\!1.6\times10^{4}$ satellites. Most of the stellar mass is in the few most massive
satellites, which have $M_\star \sim 10^{7}$--$10^{8}\,M_\odot$.

\subsection{Satellite Phase-space Coordinates} \label{sec:mock_phase}

Satellites are placed at galactocentric radii drawn from a broken power-law number-density
profile,
\begin{equation}
    n(r) \propto
    \begin{cases}
        r^{\,\beta_{\rm in}},  & r_{\rm min} \le r < r_{\rm br}, \\
        r_{\rm br}^{\,\beta_{\rm in}-\beta_{\rm out}}\, r^{\,\beta_{\rm out}},
                               & r_{\rm br} \le r \le r_{\rm max},
    \end{cases}
    \label{eq:nr}
\end{equation}
which is continuous at $r_{\rm br}$. We sample it by analytically inverting the cumulative
distribution of $dN/dr = 4\pi r^2 n(r)$. In the fiducial model,
$\beta_{\rm in} = \beta_{\rm out} = -3$ and $r_{\rm br} = r_{\rm max}$, so the profile is a
single $r^{-3}$ power law between 2 and 100\,kpc, roughly matching the density profile of
the Galactic stellar halo \citep[e.g.,][]{Deason2011}. Angular positions are isotropic on the
sphere.

Satellite velocities are isotropic in direction. Their speeds are drawn uniformly from
$[0, v_{\rm esc}(r)]$, where
\begin{equation}
    v_{\rm esc}(r) = \sqrt{2\left[\Phi_{\rm host}(r_{\rm vir}) - \Phi_{\rm host}(r)\right]}
\end{equation}
is the escape speed relative to a virial radius $r_{\rm vir} = 200$\,kpc. Every satellite is
therefore bound to the host. This prescription deliberately does not assume an equilibrium
distribution function. It spans a wide range of orbital energies and eccentricities, so the
streams have a variety of morphologies, from thin, coherent streams to shell-like and
phase-mixed debris. As an alternative, \texttt{StreamHalo} can also draw speeds from the
isotropic \citet{Hernquist1990} distribution function.

The sampled coordinates $(\mathbf{x}_i, \mathbf{v}_i)$ are the \emph{present-day}
($t = 0$) phase-space coordinates of each progenitor.

\subsection{Tidal Stream Generation} \label{sec:mock_streams}

\paragraph{Particle budget.}
We assign a total of $N_{\rm stream}$ stream particles to the satellites in proportion to
stellar mass,
\begin{equation}
    N_i = \mathrm{round}\!\left(N_{\rm stream}\,
          \frac{M_{\star,i}}{\sum_j M_{\star,j}}\right),
\end{equation}
so that each stream particle represents approximately the same amount of stellar mass. We
skip satellites with $N_i < N_{\rm min} = 100$. These faint systems would contribute
negligible light and could not be resolved as streams. Because the stellar mass is dominated
by the most massive satellites, only a small fraction of the population is retained. In the
fiducial realization, 75 satellites with $M_\star \gtrsim 2\times10^{6}\,M_\odot$ produce
streams, containing $\approx 4\times10^{4}$ particles in total.

\paragraph{Progenitor model.}
Each progenitor is modeled as a Plummer sphere with initial mass $M_h$. Its scale radius is
tied to the tidal field at the progenitor's present-day radius $r_i = |\mathbf{x}_i|$,
\begin{equation}
    a_{{\rm P},i} = \eta\, r_{t,i}, \qquad
    r_{t,i} = r_i \left[\frac{M_{h,i}}{3\,M_{\rm NFW}(<r_i)}\right]^{1/3},
    \label{eq:plummer}
\end{equation}
where $r_t$ is the \citet{King1962} tidal radius, $M_{\rm NFW}(<r)$ is the enclosed mass of
the NFW component (concentration $c = r_{\rm vir}/R_s$), and $\eta = 0.1$. Satellites that
are more massive or on wider orbits are therefore more extended. The progenitor loses mass
during the integration. Its mass decreases linearly in $\log_{10} M$ from $M_h$ at the start
of the integration to $\min(M_h, 10^{8}\,M_\odot)$ at the present day.

\paragraph{Particle spray.}
Streams are generated with the particle-spray method of \citet{Chen2025}. Each progenitor's
orbit is first integrated backward in the host potential for $T_{\rm int} = 4$\,Gyr from its
present-day coordinates. It is then integrated forward, and stream particles are released at
a constant rate from both Lagrange points along this orbit, so the leading and trailing arms
receive equal numbers of particles. At each release time we compute the Jacobi radius
\begin{equation}
    r_J = \left(\frac{G M_{\rm sat}}{\Omega^2 - \partial^2\Phi_{\rm host}/\partial r^2}\right)^{1/3},
\end{equation}
where $\Omega = |\mathbf{L}|/r^2$ is the orbital angular frequency. We evaluate the radial
curvature of the potential by automatic differentiation. The six-dimensional offsets of each
released particle, $(\Delta r/r_J,\ \phi,\ \theta,\ \Delta v/v_{\rm esc},\ \alpha,\ \beta)$,
are drawn in the satellite frame from the multivariate Gaussian that \citet{Chen2025}
calibrated on $N$-body simulations. Here $v_{\rm esc} = \sqrt{2 G M_{\rm sat}/\Delta r}$ is
the local escape speed from the progenitor. The offsets are then rotated into the host frame.
Unlike the \citet{Fardal2015} prescription, this model reproduces the offset and spread of
stream particles found in $N$-body simulations, including the dependence on the angle
between the release direction and the orbit.

After release, each particle is integrated to $t = 0$ with a leapfrog integrator in the
combined potential of the host and the moving, mass-losing Plummer progenitor. We use
$n_{\rm steps} = 8000$ steps ($\Delta t = 0.5$\,Myr). The progenitor trajectory and mass at
every release time are precomputed and shared by all particles released at that time, and
the particle integrations are vectorized with \texttt{jax.vmap}. The spray generates
$n_{\rm gen} = n_{\rm steps}\lceil N_i/n_{\rm steps}\rceil \ge N_i$ particles, from which
$N_i$ are drawn at random without replacement. Each stream particle carries the index of its
parent satellite.

\subsection{Smooth Background} \label{sec:mock_bg}

Real stellar halos also contain a phase-mixed component that is not associated with any
coherent stream. We represent it with $N_{\rm bg}$ background particles drawn from the
density profile of Equation~(\ref{eq:nr}), extended to $r_{\rm max,bg} = 200$\,kpc, with
isotropic angular positions. The number of background particles is set by the stream
fraction $f_{\rm stream}$,
\begin{equation}
    N_{\rm bg} = N_{\rm stream}^{\rm act}\,\frac{1 - f_{\rm stream}}{f_{\rm stream}},
\end{equation}
where $N_{\rm stream}^{\rm act}$ is the number of stream particles actually produced after
low-mass satellites are removed. Background particles are labeled with index $-1$. The
final mock halo combines the stream and background particles, each tagged with its parent
satellite. From it, we compute heliocentric sky coordinates as seen by an observer at
$(8, 0, 0)$\,kpc for the clustering analysis in Section~\ref{sec:2pcf}.

\begin{deluxetable}{llcl}
\tablecaption{Fiducial parameters of the mock stellar halos. \label{tab:mock_params}}
\tablehead{
\colhead{Component} & \colhead{Parameter} & \colhead{Value} & \colhead{Description}
}
\startdata
\multicolumn{4}{l}{\emph{Host potential}} \\
Disk (Miyamoto--Nagai) & $M$, $a$, $b$ & $4.77\times10^{10}\,M_\odot$, 2.6\,kpc, 0.3\,kpc & stellar disk \\
Bulge (Hernquist)      & $M$, $a$      & $5.0\times10^{9}\,M_\odot$, 1.0\,kpc   & bulge \\
Nucleus (Hernquist)    & $M$, $a$      & $1.81\times10^{9}\,M_\odot$, 0.069\,kpc & nucleus \\
Halo (NFW)             & $M$, $R_s$    & $5.54\times10^{11}\,M_\odot$, 15.6\,kpc & dark halo \\
\hline
\multicolumn{4}{l}{\emph{Satellite population}} \\
& $\alpha$ & $-1.9$ & halo mass-function slope \\
& $[M_{\rm min}, M_{\rm max}]$ & $[10^{7}, 3\times10^{10}]\,M_\odot$ & halo mass range \\
& $M_{\star,\rm tot}$ & $10^{9}\,M_\odot$ & total accreted stellar mass \\
& SHMR & \citet{Behroozi2013}, $z=0$ & stellar-to-halo mass relation \\
& $[r_{\rm min}, r_{\rm max}]$ & $[2, 100]$\,kpc & satellite radial range \\
& $\beta_{\rm in}, \beta_{\rm out}$ & $-3, -3$ & density slopes \\
& $r_{\rm vir}$ & 200\,kpc & zero point of $v_{\rm esc}$ \\
\hline
\multicolumn{4}{l}{\emph{Streams}} \\
& $N_{\rm stream}$ & $5\times10^{4}$ & target stream particles \\
& $N_{\rm min}$ & 100 & minimum particles per stream \\
& $\eta$ & 0.1 & Plummer scale / tidal radius \\
& $T_{\rm int}$ & 4\,Gyr & disruption time \\
& $\Delta t$ & 0.5\,Myr & integration time step \\
& spray model & \citet{Chen2025} & release distribution \\
\hline
\multicolumn{4}{l}{\emph{Background}} \\
& $f_{\rm stream}$ & 0.5 & stream fraction \\
& $r_{\rm max,bg}$ & 200\,kpc & background outer radius \\
\enddata
\end{deluxetable}
